\section{Κλιμακωσιμότητα Workers και Στρατηγικές Μετανάστευσης}
\label{sec:eval_scalability}

\subsection{Ανάλυση Κλιμακωσιμότητας Εργατών και Νόμοι Επιτάχυνσης}
Η κλιμακωσιμότητα της συστοιχίας εργατών (Worker Pool) αποτελεί θεμελιώδη δείκτη της ικανότητας του συστήματος να εκμεταλλεύεται παράλληλους υπολογιστικούς πόρους. Εξετάστηκε η συμπεριφορά του συστήματος για αριθμό εργατών $W \in \{1, 2, 4, 8, 16, 32\}$, συγκρίνοντας τη μετρούμενη επιτάχυνση με την ιδανική γραμμική κλίμακα, τον Νόμο του Amdahl ($S_{\text{Amdahl}} = \frac{1}{(1-p) + p/W}$) και τον Νόμο του Gustafson ($S_{\text{Gustafson}} = W - \alpha(W-1)$) με παράλληλο κλάσμα $p \approx 0.96$.

Στο Σχήμα~\ref{fig:worker_scalability} παρουσιάζεται η καμπύλη επιτάχυνσης και η εξελισσόμενη ρυθμαπόδοση (άτομα ανά δευτερόλεπτο). Στον Πίνακα~\ref{tab:worker_scalability} καταγράφονται ο χρόνος ολοκλήρωσης γενιάς, η επιτάχυνση και η παράλληλη αποδοτικότητα $\eta = \frac{S(W)}{W}$.

\begin{figure}[htbp]
    \centering
    \includegraphics[width=0.92\textwidth]{figures/fig4_worker_scalability.png}
    \caption{Κλιμακωσιμότητα συστοιχίας εργατών CFD: (α) Μετρούμενη επιτάχυνση έναντι θεωρητικών νόμων Amdahl, Gustafson και ιδανικής γραμμικής επιτάχυνσης για $W=1\dots 32$ εργάτες, και (β) ρυθμαπόδοση αξιολόγησης (άτομα/sec).}
    \label{fig:worker_scalability}
\end{figure}

../../evaluation/figures/table_fig4_worker_scalability.tex

Από τον Πίνακα~\ref{tab:worker_scalability} προκύπτει ότι το σύστημα διατηρεί εξαιρετικά υψηλή παράλληλη αποδοτικότητα $\eta > 89\%$ έως και τους 8 εργάτες (επιτάχυνση $7.15\times$). Στους 16 εργάτες η επιτάχυνση ανέρχεται σε $12.84\times$ ($\eta = 80.3\%$), ενώ στους 32 εργάτες επιτυγχάνεται ταχύτητα $21.54\times$ με ρυθμαπόδοση $98.2\text{ άτομα/s}$. Η ελαφρά πτώση της αποδοτικότητας στους 32 εργάτες οφείλεται στον κορεσμό του δικτύου gRPC και στην αναμονή συγχρονισμού τέλους γενιάς του Orchestrator.

\subsection{Ανάλυση Υπερπαραμέτρων Μετανάστευσης Νήσων}
Στην κατανεμημένη αρχιτεκτονική νησιών, η συχνότητα μετανάστευσης (διάστημα γενεών $\tau$) και το πλήθος των ανταλλασσόμενων κορυφαίων ατόμων ($M_c$) ρυθμίζουν την ισορροπία μεταξύ τοπικής εξερεύνησης (exploration) και παγκόσμιας σύγκλισης (exploitation). Αξιολογήθηκαν διαμορφώσεις με $\tau \in \{2, 5, 15\}$ και $M_c \in \{1, 3\}$.

Στο Σχήμα~\ref{fig:island_migration_tradeoffs} απεικονίζονται οι τροχιές σύγκλισης καταλληλότητας και η διατήρηση εντροπίας Shannon. Στον Πίνακα~\ref{tab:migration_tradeoffs} συνοψίζονται τα χαρακτηριστικά συμπεριφοράς κάθε διαμόρφωσης.

\begin{figure}[htbp]
    \centering
    \includegraphics[width=0.95\textwidth]{figures/fig4_island_migration_tradeoffs.png}
    \caption{Ανάλυση παραμέτρων μετανάστευσης νησιών: (α) Τροχιές εξέλιξης καταλληλότητας και (β) διατήρηση εντροπίας Shannon για διαφορετικούς συνδυασμούς διαστήματος $\tau$ και πλήθους μεταναστών $M_c$.}
    \label{fig:island_migration_tradeoffs}
\end{figure}

../../evaluation/figures/table10_distributed_migration_and_worker_scalability.tex

Η πολύ συχνή μετανάστευση ($\tau=2, M_c=3$) προκαλεί πρόωρη ομογενοποίηση των υπο-πληθυσμών, περιορίζοντας την τελική εντροπία στο $0.495$. Αντίθετα, η αραιή μετανάστευση ($\tau=15, M_c=3$) καθυστερεί τη διάχυση των καλών γονιδίων, απαιτώντας 8 γενιές για σύγκλιση. Η βέλτιστη ισορροπία επιτυγχάνεται για $\tau=5$ και $M_c=3$, όπου το σύστημα επιτυγχάνει τη μέγιστη αεροδυναμική απόδοση ($L/D = 14.64$) σε μόλις 5 γενιές, διατηρώντας ταυτόχρονα υψηλή γενετική ποικιλομορφία ($H = 0.503$).

\subsection{Πολυ-Αντικειμενική Βελτιστοποίηση και Μέτωπο Pareto}
Η σχεδίαση ενός αεροσκάφους απαιτεί τη συνδυαστική βελτιστοποίηση αντικρουόμενων αεροδυναμικών και δομικών στόχων. Μελετήθηκαν πέντε κρίσιμες παράμετροι: η αεροδυναμική απόδοση ($L/D$), η συνολική δομική μάζα ($m$), ο εσωτερικός ωφέλιμος όγκος ($V_{\text{fuse}}$), το περιθώριο στατικής διαμήκους ευστάθειας ($-C_{m_\alpha}$) και ο λόγος επιμήκυνσης ($AR$).

Στο Σχήμα~\ref{fig:aero_pareto_radar} παρουσιάζεται το μέτωπο Pareto και το διάγραμμα ακτινών (Radar Chart) για τρεις χαρακτηριστικές υποψήφιες διαμορφώσεις. Στον Πίνακα~\ref{tab:aero_pareto_comparison} παρατίθενται οι αντίστοιχες ποσοτικές τιμές.

\begin{figure}[htbp]
    \centering
    \includegraphics[width=0.92\textwidth]{figures/fig5_aero_pareto_radar.png}
    \caption{Πολυ-αντικειμενική αεροδυναμική ανάλυση Pareto: (α) Μέτωπο Pareto μεταξύ αεροδυναμικής απόδοσης ($L/D$) και δομικής μάζας, και (β) 5-αξονικό διάγραμμα ραντάρ συγκρίνοντας τη βασική διαμόρφωση, την εξειδίκευση μέγιστου $L/D$ και την ισορροπημένη ελίτ λύση.}
    \label{fig:aero_pareto_radar}
\end{figure}

../../evaluation/figures/table_fig5_aero_pareto_comparison.tex

Όπως φαίνεται στον Πίνακα~\ref{tab:aero_pareto_comparison}, η εξειδικευμένη λύση μεγιστοποίησης $L/D$ επιτυγχάνει $L/D = 13.28$ αυξάνοντας ελαφρώς τη μάζα στα $136.1\text{ g}$, ενώ η ισορροπημένη λύση (Pareto Balanced Elite) θυσιάζει ένα μικρό μέρος της αεροδυναμικής άντωσης ($L/D = 11.87$) προς όφελος σημαντικά ελαφρύτερης κατασκευής ($122.0\text{ g}$) και υψηλότερης στατικής ευστάθειας ($-C_{m_\alpha} = 1.582$).
